\documentclass[10pt,a4paper]{amsart}
\usepackage[margin=19mm]{geometry}
\usepackage{amsmath,amssymb,mathtools}
\usepackage[scr=rsfs,cal=euler]{mathalfa}
\usepackage{xfrac}
\usepackage[unicode,hidelinks,bookmarks=false]{hyperref}
\newcommand{\ie}{i.e.\ }
\newcommand{\cf}{cf.\ }
\newcommand{\resp}{respectively\ }
\newcommand{\Subsection}{\subsection}
\providecommand{\nobreakdash}{\nobreak\hbox{-}}
\setlength{\emergencystretch}{2em}
\multlinegap0pt
\usepackage{bm}
\usepackage{bbm}
\usepackage{dsfont}
\usepackage{xspace}
\usepackage{cancel}
\usepackage{mathtools}
\usepackage{amsthm}
\makeatletter
\@ifundefined{lem}{\newtheorem{lem}{Lemma}}{}
\@ifundefined{thm}{\newtheorem{thm}{Theorem}}{}
\makeatother
\title[A quotient-lifting approach to the Hamiltonicity of the cyli]{A quotient-lifting approach to the Hamiltonicity of the cylindrical 5-puzzle graph}
\author{Taizo Sadahiro}
\date{Source: arXiv:2507.03926v3 (CC BY 4.0)}
\begin{document}
\maketitle

\noindent\emph{Source and licence.} A quotient-lifting approach to the Hamiltonicity of the cylindrical 5-puzzle graph. Version arXiv:2507.03926v3 (CC BY 4.0), distributed under the Creative Commons Attribution 4.0 International licence, as recorded in the arXiv licence metadata for this exact version and shown on the abstract page https://arxiv.org/abs/2507.03926v3. Full rights notice: licenses/arxiv-2507.03926-CC-BY-4.0.txt.\par\smallskip

\noindent\textit{Editorial scope.} Counted unit: the Thm whose source label is \texttt{thm:path} in the article (source0.tex lines 389--411, source label \texttt{thm:path}), delivered together with the article's own complete original proof of it (source0.tex lines 413--494, one \texttt{proof} environment of 82 lines). No mathematics is written, repaired or re-derived: statement and proof are the article's own text line for line. Required context retained uncounted: lem:hamS5 (lines 256--286); eq:s (lines 259--263); eq:s' (lines 265--268); eq:w (lines 271--277). Every definition, display and label of those lines is retained unchanged. Omitted from this package and not redistributed: 1--255, 264--264, 269--270, 278--388, 412--412, 495--744. Nothing inside a retained excerpt is omitted or altered: every retained body is delivered exactly as the article's lines are written, so no omission receipt of any kind accompanies this package. Citations: the retained excerpts carry no citation command at all, so no bibliography entry of the article is transcribed and no reference is redistributed. Reference adaptation: none was needed and none was made; every reference inside the delivered unit resolves inside it, so no unresolved reference or citation is produced. Printed slips kept literal: no printed word, symbol, hypothesis or number of the counted statement or of its proof was altered. Layout and class: the article's own class is \texttt{article} with packages including amssymb, amsthm, float, tikz, pgf, graphicx; the wrapper is the lane's pinned \texttt{amsart} editorial wrapper at 10pt on A4, which restates the theorem-like environments the retained text needs and adds the article's own macro definitions that the retained text needs (none), with extra wrapper packages (bm, bbm, dsfont, xspace, cancel, mathtools). The delivered numbering is the unit's own. Assets: no retained span contains a figure environment, a graphics-inclusion command, a PDF-page inclusion, a table or a TikZ picture; no image file is copied into this package, the package declares no asset dependency and no image file is redistributed with it. Licence: version arXiv:2507.03926v3 of this article is distributed under the Creative Commons Attribution 4.0 International licence, as recorded in the arXiv licence metadata for this exact version and shown on the abstract page https://arxiv.org/abs/2507.03926v3. A full rights notice is delivered at licenses/arxiv-2507.03926-CC-BY-4.0.txt. Characters: every retained excerpt is pure ASCII, so the delivered engine, XeLaTeX with its default Latin Modern text font, reports no missing character.
\begin{lem}
 \label{lem:hamS5}
Let \( s = (s_1, s_2, \ldots, s_{12}) \in \{\mathtt{L}, \mathtt{R}, \mathtt{V}\}^{12} \) be the sequence
\begin{equation}
\label{eq:s}
s = (\mathtt{V}, \mathtt{L}, \mathtt{V}, \mathtt{R}, 
\mathtt{V}, \mathtt{L}, \mathtt{V}, \mathtt{R}, \mathtt{V}, \mathtt{R}, \mathtt{V}, \mathtt{L}), 
\end{equation}
and let \( s' = (s_{13}, s_{14}, \ldots, s_{24}) \) be its perturbation defined by
\begin{equation}
 \label{eq:s'}
s_{12+i} = s_i \quad \mbox{for } i = 1, 2, \ldots, 11, \qquad s_{24} = \mathtt{R}.
\end{equation}
 Define $w$ as the five-fold repetition of $(s_1,s_2,\ldots, s_{24})$,
 that is,
 \begin{equation}
  \label{eq:w}
 w  = (s_1,\ldots, s_{24}, s_1, \ldots, s_{24}, s_1, \ldots, s_{24},
 s_1, \ldots, s_{24}, s_1, \ldots, s_{24}
 )
 \in S_5^{120}.
 \end{equation}
 Then the sequence 
 \[
 %C = ({\rm id}, \phi_{S_5}(w_1), \phi_{S_5}(w_1,w_2), \ldots, \phi_{S_5}(w_1,\ldots,w_{120})) \in S_5^{121}
  C = ({\rm id}, g_1, g_2, \ldots, g_{120}) \in S_5^{121}
 \]
 is a Hamiltonian cycle in the
 Cayley graph ${\rm Cay}(S_5,\{{\mathtt L}, {\mathtt R}, {\mathtt V}\})$, where
 $g_i=\phi_{S_5}(w_1,w_2,\ldots, w_i)$.
\end{lem}

\begin{thm}
\label{thm:path}
The state graph ${\rm puz}(\Gamma_{2,3})$ of the cylindrical $5$-puzzle
has a Hamiltonian path which has the following explicit expression.
Let $s$ and $s'$ be defined by $(\ref{eq:s})$ and $(\ref{eq:s'})$
respectively, and let $t$ and $t'$ be their rotations, that is,
\[
t = (\mathtt{L}, \mathtt{V}, \mathtt{R}, 
\mathtt{V}, \mathtt{L}, \mathtt{V}, \mathtt{R}, \mathtt{V}, \mathtt{R}, \mathtt{V}, \mathtt{L},\mathtt{V})
\]
and
\[
t' = (\mathtt{L}, \mathtt{V}, \mathtt{R}, 
\mathtt{V}, \mathtt{L}, \mathtt{V}, \mathtt{R}, \mathtt{V}, \mathtt{R}, \mathtt{V}, \mathtt{R},\mathtt{V}).
\]
Then define $v$ as the word obtained by deleting the last letter from
the word
\[
 (ss')^{14}\cdot s s\cdot (tt')^{15},
\]
so that $v$ has length $719$.
Then $\hat{v}$ traces a Hamiltonian path of the state graph of the cylindrical $5$-puzzle.
\end{thm}

\begin{proof}
 Let $w$ be defined by $(\ref{eq:w})$.
 Then, as we have shown in Lemma $\ref{lem:hamS5}$,
 $\phi_{S_5}(w)={\rm id}\in S_5$.
 The word $ss'$ of length $24$ contains seven ${\mathtt R}$s and 
 five ${\mathtt L}$s.
 Therefore, the ${\mathbb Z}/3{\mathbb Z}$-part
 of $\phi_G(\hat{ss'})$ is
 \[
  5\times 1 + 7\times 2 = 1\in {\mathbb Z}/3{\mathbb Z}.
 \]
 Hence the ${\mathbb Z}/3{\mathbb Z}$-part of
 $\phi_G(\hat{w})$ is $5\times{1}=2\in {\mathbb Z}/3{\mathbb Z}$.
 In the same manner, we see that
 the ${\mathbb Z}/2{\mathbb Z}$-part of
 $\phi_G(\hat{w})$ is $0\in {\mathbb Z}/2{\mathbb Z}$.
 Define 
 \[
  d=(d_1,d_2,\ldots,d_{360})=\hat{w}^3\in G^{360}.
 \]
 and
 \[
 q_i=\phi_G(d_1\cdots d_i).
 \]
 Let
 \[
 a_0 = \left({\rm id}, 0, 0\right), ~~~
 a_1 = \left({\rm id}, 1, 0\right)\in G,
 \]
 and define
 \[
  C_i = (a_i, a_iq_1, a_iq_2, \ldots, a_iq_{360})
 \]
 for $i=0,1$. Then, now we claim that $C_1$ and $C_2$ form
 a $2$-cycle cover of the state graph ${\rm puzz}$.
 To prove this, it suffices to show that each $C_i$ is
 a simple cycle and $C_1$ and $C_2$ do not intersect.
 Suppose for the sake of contradiction, 
 that there exist $i$ and $j$ such that 
 $q_i=q_j$ and $1\leq i < j\leq 360$.
 Then, since their $S_5$-parts must coincide,
 there exists an integer $k\in\{1,2\}$ such that $j=i+120k$,
 where the ${\mathbb Z}/3{\mathbb Z}$-parts do not
 coincide. Thus $C_0$ and $C_1$ are both simple cycles.
 If we have $a_0q_i=a_1q_j$, again, since 
 since their $S_5$-parts must coincide,
 there exists an integer $k\in\{1,2\}$ such that $j=i+120k$,
 where the ${\mathbb Z}/2{\mathbb Z}$-parts do not
 coincide.

 Then we splice $C_0$ and $C_1$ to form a Hamiltonian path
 in ${\rm puz}(\Gamma_{2,3})$.
 To break the cycle of length $360$,
 we modify the last letter $d_{360}={\mathtt R}$ of $d$ to ${\mathtt L}$.
 Then
 \[
  \phi_G(d_1,d_2,\ldots,d_{359},\hat{{\mathtt L}}) =
 a_0q_{360}\hat{\mathtt R}^{-1}\hat{{\mathtt L}}
 =
 \hat{{\mathtt L}}^2 
 =(12453,0,1) = (g_1, 0, 1).
 \]
 Therefore, setting $v'=(v_1,\ldots, v_{360})=(tt')^{15}$,
 we obtain
 \[
 (
 \hat{{\mathtt L}}^2,
 \hat{{\mathtt L}}^2\phi_G(v_1),
 \hat{{\mathtt L}}^2\phi_G(v_1,v_2),
 \ldots,
 \hat{{\mathtt L}}^2\phi_G(v_1,\ldots,v_{360})
 ).
 \]
 is a simple cycle in ${\rm puz}(\Gamma_{2,3})$
 starting from and terminating
 at $\hat{{\mathtt L}}^2$, which is a cyclic shift of $C_1$
 and disjoint with $C_0$.
 Thus, 
 the word obtained by removing the last letter from the word
 $(ss')^{14}ss(tt')^{15}$ traces out a Hamiltonian path
 of ${\rm puz}(\Gamma_{2,3})$.
\end{proof}

\end{document}
